\documentclass[12pt, a4paper]{article}

% --- 宏包引入 ---
\usepackage[UTF8]{ctex}       % 中文支持
\xeCJKDeclareCharClass{CJK}{"2460->"24FF} % 让圆圈数字使用中文字体
\usepackage{amsmath, amssymb} % 数学公式
\usepackage{geometry}         % 页面设置
\usepackage{xcolor}           % 颜色支持
\usepackage{pgfplots}         % 绘图神器
\usepackage{tikz}             % 绘图基础
\usepackage[most]{tcolorbox}  % 美观的文本框
\usepackage{fancyhdr}         % 页眉页脚
\usepackage{enumitem}         % 列表格式
\usepackage{multicol}         % 多栏排版
\usepackage{siunitx}          % 单位支持

% --- 列表全局设置 ---
\setlist[enumerate]{itemsep=2pt, parsep=0pt, topsep=2pt, partopsep=0pt}
\setlist[itemize]{itemsep=2pt, parsep=0pt, topsep=2pt, partopsep=0pt}

% --- 页面设置 ---
\geometry{left=1.5cm, right=1.5cm, top=1.5cm, bottom=1.5cm}
\pgfplotsset{compat=1.18} 

% --- 配色方案 (保持与您之前的风格一致) ---
\definecolor{mainblue}{RGB}{0, 102, 204}      
\definecolor{maingreen}{RGB}{0, 128, 0}      
\definecolor{mainred}{RGB}{204, 0, 0}
\definecolor{boxbg}{RGB}{245, 247, 250}
\definecolor{highlight}{RGB}{255, 240, 200} % 高亮背景

% --- 自定义盒子 ---

% 1. 知识点盒子
\newtcolorbox{conceptbox}[1]{
    enhanced, colback=white, colframe=mainblue,
    fonttitle=\bfseries\large, title={#1},
    attach boxed title to top left={yshift*=-\tcboxedtitleheight/2, xshift=5mm},
    boxed title style={colback=mainblue, sharp corners},
    boxrule=1pt, top=10pt, bottom=5pt, drop fuzzy shadow
}

% 2. 口诀/技巧盒子 (新增)
\newtcolorbox{tipbox}{
    enhanced, colback=highlight, colframe=orange!80!red,
    coltitle=white, fonttitle=\bfseries, title={【秒杀口诀】},
    boxrule=0.5pt, leftrule=4pt, arc=2mm,
    top=4pt, bottom=4pt
}

% 3. 例题盒子
\newtcolorbox{examplebox}{
    enhanced, colback=boxbg, colframe=gray!50,
    coltitle=black, fonttitle=\bfseries, title={【精选例题】},
    detach title, before upper={\tcbtitle\par\vspace{1pt}},
    boxrule=0.5pt, leftrule=3pt, colframe=mainblue!50,
    arc=0mm, top=4pt, bottom=4pt
}

% 4. 答题区域盒子
\newtcolorbox{answerblank}[1][]{
    enhanced, colback=white, colframe=gray!40,
    boxrule=0.5pt, top=15pt, bottom=15pt,
    overlay={\node[anchor=north west, font=\small\itshape, gray] at (frame.north west) {[#1答题区域]};}
}

% --- 页眉页脚 ---
\pagestyle{fancy}
\fancyhf{}
\fancyhead[L]{\bfseries 初中数学 · 核心考点}
\fancyhead[R]{\bfseries 二次函数 · 入门与图象}
\fancyfoot[C]{\thepage}

\begin{document}

\begin{center}
    {\zihao{2} \heiti 二次函数：从零基础到中考真题} \\
    \vspace{0.1cm}
    {\large \kaishu 第一讲：图象性质与平移变换 \textbf{（学生版）}}
\end{center}

\vspace{0.1cm}

\begin{center}
    \fcolorbox{mainred}{white}{\parbox{0.9\textwidth}{
        \centering
        \textbf{姓名：\underline{\hspace{4cm}} \quad 班级：\underline{\hspace{3cm}} \quad 日期：\underline{\hspace{3cm}}}
    }}
\end{center}

\vspace{0.3cm}

% ==========================================
% 第一层级
% ==========================================
\section{第一层级：认识二次函数}
\textit{目标：能够识别二次函数，掌握各项系数的含义。}

\begin{conceptbox}{考点一：定义与一般式}
    \begin{itemize}
        \item \textbf{一般式：} $y = ax^2 + bx + c$ \quad ($a \neq 0$)
        \item \textbf{要素：} 
        \begin{enumerate}
            \item 最高次数必须为 \textbf{2}。
            \item 二次项系数 $a$ 不能为 \textbf{0} (否则就变成一次函数了)。
            \item 整式函数 (分母中不能含有自变量)。
        \end{enumerate}
    \end{itemize}
\end{conceptbox}

\begin{examplebox}
    \begin{enumerate}[label=\textbf{例\arabic*.}]
        \item 下列函数中，是二次函数的是（ \quad ）\\
        A. $y = 2x + 1$ \quad B. $y = (x-1)^2 - x^2$ \quad C. $y = \dfrac{1}{x^2} + 2$ \quad D. $y = 1 - x^2$
        
        \item \textbf{（概念陷阱）} 若函数 $y = (m-2)x^{m^2-2} + 3x$ 是关于 $x$ 的二次函数，则 $m$ 的值为 \underline{\hspace{3cm}}。
        
        \item 将二次函数 $y = 2(x-1)(x+3)$ 化为一般式 ($y=ax^2+bx+c$)，其中 $a=$\underline{\hspace{1.5cm}}, $b=$\underline{\hspace{1.5cm}}, $c=$\underline{\hspace{1.5cm}}。
    \end{enumerate}
\end{examplebox}

\begin{answerblank}[第1-3题]
    \vspace{4cm}
\end{answerblank}

\newpage

% ==========================================
% 第二层级
% ==========================================
\section{第二层级：图象的"长相"与 $a$ 的秘密}
\textit{目标：看到函数解析式，脑海中能浮现出图象的大致形状。}

\begin{conceptbox}{考点二：抛物线的开口与宽窄}
    对于 $y = ax^2$ (或一般式)：
    \begin{itemize}
        \item \textbf{开口方向：} 看 $a$ 的符号。
        \begin{itemize}
            \item $a > 0$：开口\textbf{向上} (像笑脸 $\smile$)，有\textbf{最底点} (最小值)。
            \item $a < 0$：开口\textbf{向下} (像哭脸 $\frown$)，有\textbf{最高点} (最大值)。
        \end{itemize}
        \item \textbf{开口大小：} 看 $|a|$ 的大小。
        \begin{itemize}
            \item $|a|$ 越大，开口越\textbf{小} (越"瘦")。
            \item $|a|$ 越小，开口越\textbf{大} (越"胖")。
        \end{itemize}
    \end{itemize}
\end{conceptbox}

% 绘图演示
\begin{center}
\begin{tikzpicture}[scale=0.7]
    \begin{axis}[
        width=7cm, axis lines=middle, xtick=\empty, ytick=\empty,
        title={$a > 0$ (开口向上)}, xlabel=$x$, ylabel=$y$,
        ymin=-1, ymax=5
    ]
        \addplot[domain=-2:2, smooth, thick, mainblue] {x^2};
        \node[mainblue] at (axis cs:1.5,2.5) {$y=x^2$};
        \addplot[domain=-1.5:1.5, smooth, thick, maingreen, dashed] {2*x^2};
        \node[maingreen] at (axis cs:0.5,3) {$y=2x^2$};
    \end{axis}
\end{tikzpicture}
\hspace{1cm}
\begin{tikzpicture}[scale=0.7]
    \begin{axis}[
        width=7cm, axis lines=middle, xtick=\empty, ytick=\empty,
        title={$a < 0$ (开口向下)}, xlabel=$x$, ylabel=$y$,
        ymin=-5, ymax=1
    ]
        \addplot[domain=-2:2, smooth, thick, mainred] {-x^2};
        \node[mainred] at (axis cs:1.5,-2.5) {$y=-x^2$};
    \end{axis}
\end{tikzpicture}
\end{center}

\begin{examplebox}
    \begin{enumerate}[label=\textbf{例\arabic*.}]
        \setcounter{enumi}{3}
        \item 在平面直角坐标系中，抛物线 $y = 3x^2$ 与 $y = -\dfrac{1}{3}x^2$ 的共同特点是（ \quad ）\\
        A. 开口都向上 \quad B. 对称轴都是 $y$ 轴 \quad C. 都有最高点 \quad D. 开口大小相同
        
        \item 已知点 $A(1, y_1)$， $B(2, y_2)$， $C(-3, y_3)$ 都在抛物线 $y = -x^2$ 上，则 $y_1, y_2, y_3$ 的大小关系是 \underline{\hspace{3cm}}。
    \end{enumerate}
\end{examplebox}

\begin{answerblank}[第4-5题]
    \vspace{4cm}
\end{answerblank}

\newpage

% ==========================================
% 第三层级
% ==========================================
\section{第三层级：顶点式与图象平移}
\textit{目标：中考必考！掌握"顶级"口诀，不用画图也能写出解析式。}

\begin{conceptbox}{考点三：顶点式 $y = a(x-h)^2 + k$}
    \begin{itemize}
        \item \textbf{顶点坐标：} $(h, k)$
        \item \textbf{对称轴：} 直线 $x = h$
        \item \textbf{最值：} 当 $x=h$ 时，函数取得最值 $y=k$。
    \end{itemize}
\end{conceptbox}

\begin{tipbox}
    \textbf{平移口诀：左加右减，上加下减}
    \begin{itemize}
        \item \textbf{括号内（x）：} 往左移加，往右移减。（针对 $h$）
        \item \textbf{括号外（y）：} 往上移加，往下移减。（针对 $k$）
    \end{itemize}
    \textit{秘籍：先把顶点写出来，移动顶点，再写回解析式，永远不会错！}
\end{tipbox}

\begin{examplebox}
    \begin{enumerate}[label=\textbf{例\arabic*.}]
        \setcounter{enumi}{5}
        \item 抛物线 $y = -2(x+3)^2 - 5$ 的顶点坐标是 \underline{\hspace{3cm}}，对称轴是直线 \underline{\hspace{3cm}}。
        
        \item 将抛物线 $y = 3x^2$ 向右平移 2 个单位，再向上平移 1 个单位，得到的抛物线解析式为 \underline{\hspace{4cm}}。
        
        \item \textbf{（逆向思维）} 将抛物线 $y = (x-1)^2 + 3$ 向左平移 2 个单位，再向下平移 4 个单位，得到的新抛物线与 $y$ 轴的交点坐标是 \underline{\hspace{3cm}}。
    \end{enumerate}
\end{examplebox}

\begin{answerblank}[第6-8题]
    \vspace{5cm}
\end{answerblank}

\newpage

% ==========================================
% 第四层级
% ==========================================
\section{第四层级：图象特征综合判断}
\textit{目标：中考选择题压轴基础，根据图象判断 $a, b, c$ 的符号。}

\begin{conceptbox}{考点四：看图说话 ($a, b, c$ 决定了什么)}
    \begin{itemize}
        \item \textbf{$a$：} 开口方向 (上正下负)。
        \item \textbf{$c$：} 与 $y$ 轴交点。若交在 $y$ 轴正半轴，则 $c>0$。
        \item \textbf{$b$：} 对称轴 $x = -\dfrac{b}{2a}$。
        \begin{itemize}
            \item \textbf{左同右异：} 
            对称轴在 $y$ 轴左侧 $\Rightarrow a, b$ 同号；
            对称轴在 $y$ 轴右侧 $\Rightarrow a, b$ 异号。
        \end{itemize}
    \end{itemize}
\end{conceptbox}

\begin{examplebox}
    \begin{enumerate}[label=\textbf{例\arabic*.}]
        \setcounter{enumi}{8}
        \item 如图所示是二次函数 $y = ax^2 + bx + c$ 的图象，则下列结论正确的是（ \quad ）
        
        \begin{minipage}{0.3\textwidth}
            \begin{tikzpicture}[scale=0.6]
                \begin{axis}[
                    axis lines=middle, xtick=\empty, ytick=\empty,
                    xlabel=$x$, ylabel=$y$, xmin=-2, xmax=4, ymin=-2, ymax=4
                ]
                    \addplot[domain=-1.5:3.5, smooth, thick, mainblue] {-(x-1)^2 + 3};
                    \draw[dashed] (axis cs:1,0) -- (axis cs:1,3);
                    \node[below] at (axis cs:1,0) {$1$};
                \end{axis}
            \end{tikzpicture}
        \end{minipage}
        \begin{minipage}{0.65\textwidth}
            A. $a > 0$ \\
            B. $c < 0$ \\
            C. $b < 0$ \\
            D. $b^2 - 4ac > 0$
        \end{minipage}
        
        \vspace{0.5cm}
        
        \item 二次函数 $y = ax^2 + bx + c$ 的图象如图所示，则点 $M\left(\dfrac{b}{a}, \dfrac{c}{a}\right)$ 在（ \quad ）
        
        \begin{minipage}{0.3\textwidth}
            \begin{tikzpicture}[scale=0.6]
                \begin{axis}[
                    axis lines=middle, xtick=\empty, ytick=\empty,
                    xlabel=$x$, ylabel=$y$, xmin=-2, xmax=3, ymin=-2, ymax=4
                ]
                    \addplot[domain=-1.5:2.5, smooth, thick, mainblue] {0.8*(x+0.5)^2 - 1};
                    \draw[dashed] (axis cs:-0.5,-1) -- (axis cs:-0.5,0);
                    \node[below] at (axis cs:-0.5,0) {$-1$};
                \end{axis}
            \end{tikzpicture}
        \end{minipage}
        \begin{minipage}{0.65\textwidth}
            A. 第一象限 \quad B. 第二象限 \quad C. 第三象限 \quad D. 第四象限
        \end{minipage}
        
        \vspace{0.5cm}
        
        \item 已知二次函数 $y = ax^2 + bx + c$ 的图象如图所示，下列结论：
        \begin{enumerate}
            \item[(1)] $a+b+c > 0$
            \item[(2)] $4a+2b+c > 0$ 
            \item[(3)] $2a+b = 0$
            \item[(4)] $3a+c > 0$
        \end{enumerate}
        其中正确的个数是（ \quad ）
        
        \begin{minipage}{0.3\textwidth}
            \begin{tikzpicture}[scale=0.6]
                \begin{axis}[
                    axis lines=middle, xtick=\empty, ytick=\empty,
                    xlabel=$x$, ylabel=$y$, xmin=-2, xmax=4, ymin=-2, ymax=4
                ]
                    \addplot[domain=-1.5:3.5, smooth, thick, mainblue] {-0.8*(x-1)^2 + 2.5};
                    \draw[dashed] (axis cs:1,0) -- (axis cs:1,2.5);
                    \node[below] at (axis cs:1,0) {$1$};
                    \fill (axis cs:-0.5,0) circle (2pt) node[below] {$-1$};
                    \fill (axis cs:2.5,0) circle (2pt) node[below] {$2$};
                \end{axis}
            \end{tikzpicture}
        \end{minipage}
        \begin{minipage}{0.65\textwidth}
            A. 1个 \quad B. 2个 \quad C. 3个 \quad D. 4个
        \end{minipage}
        
        \vspace{0.5cm}
        
        \item 如图，抛物线 $y = ax^2 + bx + c$ 经过点 $(-1,0)$，对称轴为直线 $x=1$。下列结论：
        \begin{enumerate}
            \item[(1)] $abc < 0$
            \item[(2)] $b < c$
            \item[(3)] $3a+c = 0$
            \item[(4)] $a+b \leqslant m(am+b)$（$m$为任意实数）
        \end{enumerate}
        其中正确的有（ \quad ）
        
        \begin{minipage}{0.3\textwidth}
            \begin{tikzpicture}[scale=0.6]
                \begin{axis}[
                    axis lines=middle, xtick=\empty, ytick=\empty,
                    xlabel=$x$, ylabel=$y$, xmin=-2, xmax=4, ymin=-4, ymax=2
                ]
                    \addplot[domain=-1.5:3.5, smooth, thick, mainblue] {0.5*(x-1)^2 - 2};
                    \draw[dashed] (axis cs:1,-2) -- (axis cs:1,0);
                    \node[below] at (axis cs:1,0) {$1$};
                    \fill (axis cs:-1,0) circle (2pt) node[above left] {$(-1,0)$};
                    \fill (axis cs:3,0) circle (2pt) node[above right] {$(3,0)$};
                \end{axis}
            \end{tikzpicture}
        \end{minipage}
        \begin{minipage}{0.65\textwidth}
            A. 1个 \quad B. 2个 \quad C. 3个 \quad D. 4个
        \end{minipage}
    \end{enumerate}
\end{examplebox}

\begin{answerblank}[第9-12题]
    \vspace{8cm}
\end{answerblank}

\newpage

% ==========================================
% 第五层级
% ==========================================
\section{第五层级：解一元二次方程的四大方法}
\textit{目标：熟练掌握各种解法，能根据题目特点选择最优解法。}

\begin{conceptbox}{考点五：四种解法速览}
    \begin{enumerate}
        \item \textbf{直接开平方法}：适用于 $(x+m)^2 = n$ 形式
        \item \textbf{配方法}：通用方法，关键步骤：移项 $\to$ 配方 $\to$ 开方
        \item \textbf{公式法}：万能方法，求根公式 $x = \dfrac{-b \pm \sqrt{b^2-4ac}}{2a}$
        \item \textbf{因式分解法}：最快方法，适用于能分解的题目
    \end{enumerate}
\end{conceptbox}

\begin{tipbox}
    \textbf{解法选择口诀：}
    \begin{itemize}
        \item 平方形式直接开（如 $(x-2)^2 = 9$）
        \item 能分就分最省事（如 $x^2 - 5x + 6 = 0$）
        \item 配方公式是备胎（什么题都能做）
    \end{itemize}
\end{tipbox}

\begin{examplebox}
    \begin{enumerate}[label=\textbf{例\arabic*.}]
        \setcounter{enumi}{12}
        \item 用适当的方法解方程：
        \begin{enumerate}
            \item[(1)] $(2x-1)^2 = 16$
            \item[(2)] $x^2 - 6x + 5 = 0$
            \item[(3)] $2x^2 - 4x - 1 = 0$
        \end{enumerate}
        
        \item 若关于 $x$ 的方程 $x^2 - 4x + m = 0$ 能用因式分解法求解，则 $m$ 的值可能是（ \quad ）\\
        A. 3 \quad B. 5 \quad C. 6 \quad D. 8
    \end{enumerate}
\end{examplebox}

\begin{answerblank}[第13-14题]
    \vspace{6cm}
\end{answerblank}

\newpage

% ==========================================
% 第六层级
% ==========================================
\section{第六层级：判别式——根的"体检报告"}
\textit{目标：通过判别式判断方程根的情况，中考必考！}

\begin{conceptbox}{考点六：判别式 $\Delta = b^2 - 4ac$}
    \begin{itemize}
        \item $\Delta > 0$：方程有 \textbf{两个不相等} 的实数根
        \item $\Delta = 0$：方程有 \textbf{两个相等} 的实数根（一个重根）
        \item $\Delta < 0$：方程 \textbf{没有} 实数根（在实数范围内无解）
    \end{itemize}
    \textit{注意：使用判别式前，务必先将方程化为一般式 $ax^2 + bx + c = 0$！}
\end{conceptbox}

\begin{examplebox}
    \begin{enumerate}[label=\textbf{例\arabic*.}]
        \setcounter{enumi}{14}
        \item 不解方程，判断下列方程根的情况：
        \begin{enumerate}
            \item[(1)] $2x^2 - 3x + 1 = 0$
            \item[(2)] $x^2 - 2x + 3 = 0$
            \item[(3)] $x^2 - 4x + 4 = 0$
        \end{enumerate}
        
        \item 若关于 $x$ 的一元二次方程 $kx^2 - 2x - 1 = 0$ 有两个不相等的实数根，则 $k$ 的取值范围是（ \quad ）\\
        A. $k > -1$ \quad B. $k > -1$ 且 $k \neq 0$ \quad C. $k < 1$ \quad D. $k < 1$ 且 $k \neq 0$
        
        \item 已知关于 $x$ 的方程 $x^2 - (2k+1)x + k^2 = 0$ 有两个相等的实数根，求 $k$ 的值。
    \end{enumerate}
\end{examplebox}

\begin{answerblank}[第15-17题]
    \vspace{7cm}
\end{answerblank}

\newpage

% ==========================================
% 第七层级
% ==========================================
\section{第七层级：韦达定理——根与系数的"暗号"}
\textit{目标：不解方程，也能知道两根之和与两根之积。}

\begin{conceptbox}{考点七：根与系数的关系}
    若 $x_1, x_2$ 是方程 $ax^2 + bx + c = 0$（$a \neq 0$）的两个根，则：
    \begin{itemize}
        \item \textbf{两根之和：}$x_1 + x_2 = -\dfrac{b}{a}$
        \item \textbf{两根之积：}$x_1 \cdot x_2 = \dfrac{c}{a}$
    \end{itemize}
    \textit{使用前提：方程必须有实数根（即 $\Delta \geq 0$）}
\end{conceptbox}

\begin{tipbox}
    \textbf{韦达定理常见变形：}
    \begin{itemize}
        \item $\dfrac{1}{x_1} + \dfrac{1}{x_2} = \dfrac{x_1 + x_2}{x_1 x_2}$
        \item $x_1^2 + x_2^2 = (x_1 + x_2)^2 - 2x_1 x_2$
        \item $(x_1 - x_2)^2 = (x_1 + x_2)^2 - 4x_1 x_2$
    \end{itemize}
\end{tipbox}

\begin{examplebox}
    \begin{enumerate}[label=\textbf{例\arabic*.}]
        \setcounter{enumi}{17}
        \item 已知方程 $2x^2 - 5x + 1 = 0$ 的两根为 $x_1, x_2$，求：
        \begin{enumerate}
            \item[(1)] $x_1 + x_2$ 和 $x_1 \cdot x_2$
            \item[(2)] $\dfrac{1}{x_1} + \dfrac{1}{x_2}$
            \item[(3)] $x_1^2 + x_2^2$
        \end{enumerate}
        
        \item 已知关于 $x$ 的方程 $x^2 - 3x + m = 0$ 的一个根是 1，求另一个根和 $m$ 的值。
        
        \item 已知 $x_1, x_2$ 是方程 $x^2 - 4x + 2 = 0$ 的两根，求以 $(x_1+1)$ 和 $(x_2+1)$ 为根的一元二次方程。
    \end{enumerate}
\end{examplebox}

\begin{answerblank}[第18-20题]
    \vspace{7cm}
\end{answerblank}

\newpage

% ==========================================
% 第八层级
% ==========================================
\section{第八层级：应用题实战（增长率与利润）}
\textit{目标：把数学知识用到生活中去，中考大题高频考点！}

\begin{conceptbox}{考点八：常见应用题模型}
    \begin{enumerate}
        \item \textbf{增长率模型：}$a(1+x)^n = b$
        \begin{itemize}
            \item $a$：原始量，$x$：增长率，$n$：期数，$b$：最终量
        \end{itemize}
        \item \textbf{利润模型：}利润 = (售价 - 进价) $\times$ 销量
        \begin{itemize}
            \item 常结合"降价促销，销量增加"的线性关系
        \end{itemize}
        \item \textbf{面积模型：}平移法求种植面积
    \end{enumerate}
\end{conceptbox}

\begin{examplebox}
    \begin{enumerate}[label=\textbf{例\arabic*.}]
        \setcounter{enumi}{20}
        \item \textbf{（增长率）}某工厂今年一月份生产零件 10 万个，第一季度共生产零件 33.1 万个。若每月的增长率相同，求这个增长率。
        
        \vspace{0.3cm}
        
        \item \textbf{（利润问题）}某商店购进一批单价为 16 元的日用品，若按每件 20 元的价格销售，每月能卖出 360 件；若按每件 25 元的价格销售，每月能卖出 210 件。假定每月销售件数 $y$（件）是价格 $x$（元/件）的一次函数。
        \begin{enumerate}
            \item[(1)] 求 $y$ 与 $x$ 之间的函数关系式；
            \item[(2)] 在商品不积压且不考虑其他因素的条件下，销售价格定为多少时，才能使每月获得最大利润？最大利润是多少？
        \end{enumerate}
        
        \vspace{0.3cm}
        
        \item \textbf{（面积问题）}如图，要利用一面墙（墙长 25 米）建羊圈，用 100 米的围栏围成总面积为 400 平方米的三个大小相同的矩形羊圈，求羊圈的边长 $AB$、$BC$ 各为多少米？
        
        \begin{center}
        \begin{tikzpicture}[scale=0.08]
            \draw[thick] (0,0) rectangle (60,30);
            \draw[thick] (20,0) -- (20,30);
            \draw[thick] (40,0) -- (40,30);
            \draw[thick, dashed] (0,0) -- (-10,0);
            \draw[thick, dashed] (0,30) -- (-10,30);
            \node[left] at (-5,15) {墙};
            \node[below] at (10,0) {$AB$};
            \node[below] at (30,-5) {$BC$};
        \end{tikzpicture}
        \end{center}
    \end{enumerate}
\end{examplebox}

\begin{answerblank}[第21-23题]
    \vspace{10cm}
\end{answerblank}

\vfill
\begin{center}
    \small \textcolor{gray}{—— 第一讲 学生版 完 ——}
\end{center}

\end{document}
